\subsection{Fundamentos de gestión, seguimiento y gestión documental} 
\label{subsec:fundamentos_gestion}

Para sustentar formalmente la propuesta, es indispensable comprender tres conceptos fundamentales que respaldan el uso de tecnologías en la educación superior:

\subsubsection{Sistemas de gestión académica en las universidades}
Un sistema de gestión académica es un programa informático diseñado para recopilar, procesar, organizar y consultar la información que se genera durante las actividades escolares de una institución \parencite{laudon2020management}. 

Dentro de las universidades existen principalmente dos tipos de sistemas:
\begin{itemize}
    \item \textbf{Sistemas de control escolar general (SIS / ERP universitario):} Son programas pensados para administrar las inscripciones, cobros de colegiaturas, actas oficiales y el historial de materias aprobadas de todos los estudiantes (como el sistema Banner o los portales de autoservicio). Aunque son indispensables para la administración general, no cuentan con herramientas para apoyar el trabajo académico que ocurre día con día dentro de una investigación.
    \item \textbf{Sistemas de gestión tutorial y seguimiento especializado:} Son aplicaciones orientadas a supervisar etapas formativas específicas, como el desarrollo de una tesis o las prácticas profesionales. Su objetivo no es cobrar ni registrar materias, sino organizar la entrega de avances, registrar compromisos entre profesores y estudiantes y medir los tiempos de respuesta para asegurar que los proyectos concluyan exitosamente. La página web que se presenta en este trabajo pertenece a este segundo tipo.
\end{itemize}

\subsubsection{Seguimiento del aprendizaje y acompañamiento tutorial en el posgrado}
A diferencia de los cursos de licenciatura, donde el aprendizaje se evalúa principalmente con exámenes y tareas cortas, en el posgrado el aprendizaje se basa en la investigación autodirigida y en la elaboración continua de un texto científico \parencite{lee2012successful}.

En este contexto, el acompañamiento tutorial representa la guía académica continua que brinda el profesor para orientar al tesista, resolver dudas metodológicas y revisar la redacción de sus capítulos \parencite{bitchener2014supervising}. Diversos autores señalan que el seguimiento del aprendizaje no debe depender únicamente de la memoria o de conversaciones informales, sino que requiere:
\begin{itemize}
    \item \textbf{Retroalimentación formativa y oportuna:} El alumno necesita saber con claridad qué partes de su borrador están bien argumentadas y qué aspectos debe corregir antes de la siguiente entrega.
    \item \textbf{Establecimiento de metas y compromisos alcanzables:} Cada reunión de asesoría debe culminar con tareas específicas fechadas, lo que permite que el estudiante autorregule su tiempo y evite bloqueos en la investigación.
    \item \textbf{Supervisión activa del progreso:} Cuando la facultad supervisa con regularidad los avances, es posible detectar a tiempo si un estudiante está perdiendo el ritmo de trabajo o enfrenta problemas personales, evitando que caiga en el rezago escolar o abandone el programa \parencite{lando2025digital}.
\end{itemize}

\subsubsection{Gestión documental académica y ciclo de vida de los archivos}
La gestión documental consiste en el conjunto de reglas, técnicas y prácticas utilizadas para administrar el flujo de los documentos que produce una organización, desde el momento en que se crean hasta su conservación final o descarte \parencite{iso15489}. 

En la redacción de una tesis de posgrado, el documento atraviesa un ciclo de vida compuesto por cuatro fases:
\begin{enumerate}
    \item \textbf{Creación y carga:} El estudiante redacta su avance y lo entrega en la plataforma escolar.
    \item \textbf{Revisión y retroalimentación:} El asesor descarga el archivo, evalúa el contenido y escribe sus comentarios de corrección.
    \item \textbf{Evolución y versionado inmutable:} El estudiante atiende las observaciones y sube un nuevo borrador. Para garantizar la trazabilidad del trabajo, la versión anterior no se sobreescribe ni se borra, sino que se resguarda de manera íntegra con un número correlativo (Versión 1, Versión 2, etc.), permitiendo que ambas partes puedan comprobar qué cambios se hicieron a lo largo del tiempo.
    \item \textbf{Aprobación y dictamen:} Una vez que el documento cumple con la calidad académica requerida, el asesor y el comité autorizan la versión final para el examen de grado.
\end{enumerate}

Para asegurar que los documentos mantengan su integridad y autenticidad durante este ciclo, la página web exige el uso de formatos no alterables como el estándar PDF y aplica reglas estrictas que impiden la eliminación accidental de archivos, garantizando que el expediente de titulación se mantenga completo y verificable.